\begin{appendix}
\chapter{Appendices}\label{chap:ann1}

\section{Appendix A: HFL System Architecture}
\label{anx:arquitectura}

The deployed architecture consists of three layers:

\begin{itemize}
    \item \textbf{Edge Layer:} 2 $\times$ ESP32-S3 (normal node + simulated node), running TinyML inference and MQTT communication encrypted with ASCON-128.
    \item \textbf{Fog Layer:} 1 $\times$ Raspberry Pi 4 (4\,GB RAM), acting as MQTT broker, local training gateway, and HTTP bridge to the server.
    \item \textbf{Cloud Layer:} 1 $\times$ PC (FastAPI server), coordinating FedAvg aggregation and the analytics dashboard.
\end{itemize}

\section{Appendix B: System Files}
\label{anx:archivos}

\begin{table}[H]
\centering
\caption{Main HFL v7 system files.}
\small
\begin{tabular}{l l l}
\toprule
\textbf{File} & \textbf{Device} & \textbf{Function} \\
\midrule
\texttt{main\_edge\_node\_normal.cpp}    & ESP32-S3 & Normal traffic node \\
\texttt{main\_edge\_node\_simulated.cpp} & ESP32-S3 & Mixed traffic node \\
\texttt{ascon128.h}                      & ESP32-S3 & ASCON-128 encryption (C++) \\
\texttt{gateway\_hfl.py}                 & Raspberry Pi 4 & Federated gateway \\
\texttt{ascon128.py}                     & RPi / PC & ASCON-128 encryption (Python) \\
\texttt{ascon\_metrics.py}               & RPi / PC & Encryption metrics \\
\texttt{server\_hfl.py}                  & PC & Central FedAvg server \\
\bottomrule
\end{tabular}
\end{table}

\section{Appendix C: Project Repository}
\label{anx:repositorio}

The complete source code, datasets, and training scripts are available in the project repository. The \texttt{main} branch contains the system with ASCON-128 encryption enabled, and the \texttt{hfl\_v7-no-ascon} branch contains the version without encryption for baseline measurement.

\section{Appendix D: Classical ML Model Configuration}
\label{anx:ml_classic}

Table~\ref{tab:ml_config} details the hyperparameter configurations selected after optimization for each classical ML model.

\begin{table}[H]
\centering
\caption{Optimized hyperparameters for classical ML models.}
\small
\begin{tabular}{l l}
\toprule
\textbf{Model} & \textbf{Optimized Hyperparameters} \\
\midrule
Decision Tree & \texttt{max\_depth=20, min\_samples\_split=5, criterion='gini'} \\
              & \texttt{class\_weight='balanced'} \\
\addlinespace
Random Forest & \texttt{n\_estimators=200, max\_depth=30, max\_features='sqrt'} \\
              & \texttt{class\_weight='balanced\_subsample'} \\
\addlinespace
Logistic Regr. & \texttt{C=1.0, penalty='l2', solver='lbfgs'} \\
               & \texttt{class\_weight='balanced', max\_iter=1000} \\
\addlinespace
SVM (Linear)  & \texttt{C=1.0, kernel='linear'} \\
              & \texttt{class\_weight='balanced'} \\
\addlinespace
SVM (RBF)     & \texttt{C=10, gamma='scale', kernel='rbf'} \\
              & \texttt{class\_weight='balanced', subsample=25000/class} \\
\bottomrule
\end{tabular}
\label{tab:ml_config}
\end{table}

\section{Appendix E: Feature Extraction Details}
\label{anx:features}

The 13 features extracted per network flow are computed in real time on the ESP32-S3 as each packet is captured. The implementation uses a sliding window approach where statistics (mean, standard deviation, min, max) are incrementally updated to avoid storing all individual packet data in memory:

\begin{itemize}
    \item \textbf{Packet count (\texttt{num\_pkts}):} Simple counter incremented per packet.
    \item \textbf{IAT statistics (\texttt{mean\_iat}, \texttt{std\_iat}, \texttt{min\_iat}, \texttt{max\_iat}):} Computed using Welford's online algorithm for numerical stability.
    \item \textbf{Packet length statistics:} Mean, standard deviation, min, and max computed incrementally.
    \item \textbf{TCP flag counters (\texttt{num\_psh}, \texttt{num\_rst}, \texttt{num\_urg}):} Direct extraction from TCP header fields.
    \item \textbf{Total bytes (\texttt{num\_bytes}):} Running sum of packet lengths.
\end{itemize}

All features are normalized using pre-computed StandardScaler parameters ($\mu_i$, $\sigma_i$) embedded as constants in the ESP32 firmware, computed from the training dataset.

\end{appendix}
